\section{結論}\label{sec:conclusion}

本研究では、密集区間マイニングの出力を点過程・生存分析の枠組みと接続し、
記述的パターンマイニングから予測的パターンマイニングへの橋渡しを行った。
提案した3つの概念——密集区間発生過程（DIOP）、密集区間間時間（IDIT）、
密集持続時間予測——は、密集区間の時間的構造を確率的にモデル化する
新たな視点を提供する。

合成データによる実験の結果、以下の知見が得られた：
(1) 間接的自己励起データにおいてHawkes-Denseモデルは
    Poisson基線と同等の予測精度を示し、$\hat{\alpha} \approx 0$ に退化すること。
    これは、データに明示的な自己励起構造がない場合にモデルが
    適切にPoisson過程に退化する（過適合しない）ことを示す。
(2) 真に自己励起的な生成過程においては、
    Hawkes-Denseモデルが自己励起パラメータを正確に回復し、
    尤度比検定で有意な改善を示すこと。
(3) IDIT分布の選択はデータの生成過程を反映し、
    変動係数（CV）が発生パターンの規則性の有用な指標となること。
(4) Weibull生存モデルは持続時間の点予測に加えて
    ハザード関数による解釈可能な特性記述を提供すること。

\subsection{限界と今後の課題}

本研究にはいくつかの限界がある。
\textbf{第一に}、実験が合成データに限定されており、
実データでの検証が最も重要な今後の課題である。
小売トランザクション（Dunnhumby等）やソーシャルメディアストリームへの
適用を具体的に計画している。
\textbf{第二に}、Hawkes過程の自己励起構造は
密集区間の発生が連鎖的な場合にのみ有効であり、
外部要因駆動型の密集区間に対しては
共変量を組み込んだ条件付き強度モデルへの拡張が必要である。
\textbf{第三に}、提案モデルは各アイテムセットを独立に扱っており、
アイテムセット間の相互作用は考慮していない。
多次元Hawkes過程~\cite{zhou2013learning}の導入により、
アイテムセット間の因果関係をモデル化することが可能と考えられる。
\textbf{第四に}、ニューラル点過程~\cite{du2016recurrent,mei2017neural,zuo2020transformer}
との性能比較を行っていない。
大規模実データにおいてはニューラルモデルが優位となる可能性があり、
Hawkes-Denseモデルとの使い分け基準の確立が求められる。

今後は、(1) 大規模実データへの適用と実証、
(2) 共変量付き条件付き強度モデルの開発、
(3) ニューラル点過程との比較実験、
(4) オンライン学習への拡張を予定している。
